% Generated by scripts in this repository from the committed
% corpora and models. See results/TABLES.md for provenance.

\begin{table}[htbp]
\centering
\caption{Descriptive information about the corpora.}
\label{tab:corpora}
\begin{tabular}{|l|cc|cc|cc|c|c|}
\hline
 & \multicolumn{2}{c|}{\textbf{French}} & \multicolumn{2}{c|}{\textbf{FIN}} & \multicolumn{2}{c|}{\textbf{FIN random}} & & \\
\textbf{Orthographic Measure} & M & SD & M & SD & M & SD & t & p \\
\hline
Neighbour words \# & 0.40 & 0.78 & 0.67 & 0.88 & 0.00 & 0.00 & -10.25 & <.001 \\
Levenshtein distance 20, mean & 3.12 & 0.48 & 3.05 & 0.45 & 4.66 & 0.14 & 5.11 & <.001 \\
Levenshtein distance 20, sd & 0.54 & 0.16 & 0.59 & 0.17 & 0.49 & 0.09 & -10.99 & <.001 \\
Spread (\#letters yielding a neighbour) & 0.34 & 0.59 & 0.58 & 0.74 & 0.00 & 0.00 & -11.62 & <.001 \\
Uniqueness point & 4.34 & 1.04 & 5.07 & 1.31 & 3.15 & 0.45 & -19.76 & <.001 \\
\hline
\end{tabular}

\vspace{2mm}
\begin{minipage}{0.9\textwidth}
\footnotesize \textbf{Note:} N = 1985 in French corpus, 2000 in both Finnish corpora. The statistical test results are between the French and the non-random Finnish corpus (Welch's $t$-test).
\end{minipage}
\end{table}


\begin{table}[htbp]
\centering
\caption{Model performance.}
\label{tab:performance}
\begin{tabular}{|l|ccc|ccc|ccc|}
\hline
 & \multicolumn{3}{c|}{\textbf{French Model}} & \multicolumn{3}{c|}{\textbf{Finnish Model}} & \multicolumn{3}{c|}{\textbf{Random Finnish Model}} \\
\textbf{Corpus size} & \multicolumn{3}{c|}{1985 items} & \multicolumn{3}{c|}{2000 items} & \multicolumn{3}{c|}{2000 items} \\
\hline
\textbf{Test Condition} & Accur.\ (\%) & \# & Errors(\#) & Accur.\ (\%) & \# & Errors(\#) & Accur.\ (\%) & \# & Errors(\#) \\
\hline
Real Words & 100.0\% & 13895 & 0 & 100.0\% & 14000 & 0 & 100.0\% & 14000 & 0 \\
Random String & 100.0\% & 1000 & 0 & 100.0\% & 1000 & 0 & 100.0\% & 1000 & 0 \\
Single Repeated Letter & 100.0\% & 1000 & 0 & 100.0\% & 1000 & 0 & 100.0\% & 1000 & 0 \\
Double Letter Substitution & 99.7\% & 1985 & 5 & 99.7\% & 2000 & 6 & 100.0\% & 2000 & 0 \\
Letter Transposition & 87.5\% & 1985 & 248 & 67.0\% & 2000 & 659 & 98.5\% & 2000 & 29 \\
\hline
\multicolumn{10}{|c|}{\textbf{Relative Positioning Priming}} \\
\hline
\textbf{Condition} & Primed (\%) & \# & Primed(\#) & Primed (\%) & \# & Primed(\#) & Primed (\%) & \# & Primed(\#) \\
\hline
SILE-SILENCE & 1.21\% & 1985 & 24 & 0.75\% & 2000 & 15 & 0.00\% & 2000 & 0 \\
SLNE-SILENCE & 0.10\% & 1985 & 2 & 0.20\% & 2000 & 4 & 0.00\% & 2000 & 0 \\
\hline
\multicolumn{10}{|c|}{\textbf{Transposed Letter Priming}} \\
\hline
\textbf{Condition} & Primed (\%) & \# & Primed(\#) & Primed (\%) & \# & Primed(\#) & Primed (\%) & \# & Primed(\#) \\
\hline
SILNECE-SILENCE & 30.43\% & 1985 & 604 & 25.95\% & 2000 & 519 & 18.95\% & 2000 & 379 \\
SILOPCE-SILENCE & 4.94\% & 1985 & 98 & 6.75\% & 2000 & 135 & 0.05\% & 2000 & 1 \\
\hline
\end{tabular}

\vspace{2mm}
\begin{minipage}{0.95\textwidth}
\footnotesize \textbf{Note:} Recognition and nonword conditions report accuracy at the 0.9 criterion, where an error is a nonword whose winning lexical unit reached threshold. The priming blocks report the proportion of primes that activated \emph{their own target's} unit above 0.5, following Dandurand et al.\ (2013, \S3.2); these are effects to be measured, not errors, and the prediction is SILE $>$ SLNE and SILNECE $>$ SILOPCE.
\end{minipage}
\end{table}
